% GENERATED FILE. Edit canonical lessons, then run npm run generate:books.
% canonical-source-hash: fnv1a64:fc07a80cc0cd4c8b
% canonical-lessons: ES-C753-musico, ES-C753-tararear, ES-C753-cancion, ES-C753-ritmo, ES-C753-componer

\chapter{Música y canciones}
\label{ch:es-musica-y-canciones}
\hlchaptermodality{753}

\begin{chapteropening}
\textbf{By the end of this chapter:} \emph{I can talk about musicians, humming a tune, songs and rhythm.}

Complete the last lesson of chapter 753: I can talk about musicians, humming a tune, songs and rhythm.
\end{chapteropening}

\section[el músico]{el músico — the musician}

\label{lesson:ES-C753-musico}



\begin{quote}
Before the new one: say the Spanish for magic, then the Spanish for to bet.
\end{quote}

\subsection*{You'll want to know: el músico}
\textbf{el músico} — \textquotedblleft{}the musician\textquotedblright{}.

\begin{quote}
Mi vecino es músico de jazz.
\end{quote}

\begin{quote}
\emph{My neighbour is a jazz musician.}
\end{quote}

\emph{La música} is music and also a woman musician.

\begin{cousinweb}
\emph{músico} comes from Latin \emph{mūsicus} \textquotedblleft{}musician\textquotedblright{}, from Greek \emph{mousikós} \textquotedblleft{}of the Muses\textquotedblright{} (EtymDB).
\end{cousinweb}

\begin{grammarlens}[title={What you've built}]
One more word for telling, explaining and describing at B1.
\end{grammarlens}

\subsection*{Your turn}
\begin{itemize}
\raggedright
  \item \emph{Say it:} \emph{el músico}
  \item \emph{Say it:} \emph{el músico}, once more
  \item \emph{Say it:} a sentence of your own with \emph{el músico}
  \item \emph{Recall:} say the Spanish for magic, then the Spanish for to bet, then say \emph{el músico} again
\end{itemize}

\begin{tcolorbox}[breakable,colback=teal!4,colframe=teal!35!black,title={Before you move on}]
What does el músico mean? (\textquotedblleft{}The musician\textquotedblright{}.) Say it once more.
\end{tcolorbox}
\section[tararear]{tararear — to hum}

\label{lesson:ES-C753-tararear}



\begin{quote}
Before the new one: say the Spanish for to bet, then the Spanish for the musician.
\end{quote}

\subsection*{You'll want to know: tararear}
\textbf{tararear} — \textquotedblleft{}to hum\textquotedblright{}.

\begin{quote}
Siempre tarareo en la ducha.
\end{quote}

\begin{quote}
\emph{I always hum in the shower.}
\end{quote}

To sing a tune without the words.

\begin{cousinweb}
\emph{tararear} comes from the syllables \emph{ta-ra-ra} people sing when they do not know the words; it is an imitative word.
\end{cousinweb}

\begin{grammarlens}[title={What you've built}]
One more word for telling, explaining and describing at B1.
\end{grammarlens}

\subsection*{Your turn}
\begin{itemize}
\raggedright
  \item \emph{Say it:} \emph{tararear}
  \item \emph{Say it:} \emph{tararear}, once more
  \item \emph{Say it:} a sentence of your own with \emph{tararear}
  \item \emph{Recall:} say the Spanish for to bet, then the Spanish for the musician, then say \emph{tararear} again
\end{itemize}

\begin{tcolorbox}[breakable,colback=teal!4,colframe=teal!35!black,title={Before you move on}]
What does tararear mean? (\textquotedblleft{}To hum\textquotedblright{}.) Say it once more.
\end{tcolorbox}
\section[la canción]{la canción — the song}

\label{lesson:ES-C753-cancion}



\begin{quote}
Before the new one: say the Spanish for the musician, then the Spanish for to hum.
\end{quote}

\subsection*{You'll want to know: la canción}
\textbf{la canción} — \textquotedblleft{}the song\textquotedblright{}.

\begin{quote}
Esta canción me recuerda a ti.
\end{quote}

\begin{quote}
\emph{This song reminds me of you.}
\end{quote}

\begin{cousinweb}
\emph{canción} is inherited from Latin \emph{cantiō}, \emph{cantiōnem} \textquotedblleft{}song, singing\textquotedblright{} (EtymDB), built on \emph{canere} \textquotedblleft{}to sing\textquotedblright{}.
\end{cousinweb}

\begin{grammarlens}[title={What you've built}]
One more word for telling, explaining and describing at B1.
\end{grammarlens}

\subsection*{Your turn}
\begin{itemize}
\raggedright
  \item \emph{Say it:} \emph{la canción}
  \item \emph{Say it:} \emph{la canción}, once more
  \item \emph{Say it:} a sentence of your own with \emph{la canción}
  \item \emph{Recall:} say the Spanish for the musician, then the Spanish for to hum, then say \emph{la canción} again
\end{itemize}

\begin{tcolorbox}[breakable,colback=teal!4,colframe=teal!35!black,title={Before you move on}]
What does la canción mean? (\textquotedblleft{}The song\textquotedblright{}.) Say it once more.
\end{tcolorbox}
\section[el ritmo]{el ritmo — rhythm; pace}

\label{lesson:ES-C753-ritmo}



\begin{quote}
Before the new one: say the Spanish for to hum, then the Spanish for the song.
\end{quote}

\subsection*{You'll want to know: el ritmo}
\textbf{el ritmo} — \textquotedblleft{}rhythm; pace\textquotedblright{}.

\begin{quote}
Esta canción tiene buen ritmo.
\end{quote}

\begin{quote}
\emph{This song has a good beat.}
\end{quote}

Also pace: \emph{el ritmo de vida} — the pace of life.

\begin{cousinweb}
\emph{ritmo} comes from Latin \emph{rhythmus}, borrowed from Greek \emph{rhythmós} \textquotedblleft{}measured flow\textquotedblright{}, from \emph{rheîn} \textquotedblleft{}to flow\textquotedblright{} (EtymDB).
\end{cousinweb}

\begin{grammarlens}[title={What you've built}]
One more word for telling, explaining and describing at B1.
\end{grammarlens}

\subsection*{Your turn}
\begin{itemize}
\raggedright
  \item \emph{Say it:} \emph{el ritmo}
  \item \emph{Say it:} \emph{el ritmo}, once more
  \item \emph{Say it:} a sentence of your own with \emph{el ritmo}
  \item \emph{Recall:} say the Spanish for to hum, then the Spanish for the song, then say \emph{el ritmo} again
\end{itemize}

\begin{tcolorbox}[breakable,colback=teal!4,colframe=teal!35!black,title={Before you move on}]
What does el ritmo mean? (\textquotedblleft{}Rhythm; pace\textquotedblright{}.) Say it once more.
\end{tcolorbox}
\section[componer]{componer — to compose, to write (music)}

\label{lesson:ES-C753-componer}



\begin{quote}
Before the new one: say the Spanish for the song, then the Spanish for rhythm.
\end{quote}

\subsection*{You'll want to know: componer}
\textbf{componer} — \textquotedblleft{}to compose, to write (music)\textquotedblright{}.

\begin{quote}
Compone canciones para su grupo.
\end{quote}

\begin{quote}
\emph{He writes songs for his band.}
\end{quote}

Conjugates like \emph{poner}: \emph{compongo}, \emph{compuse}.

\begin{cousinweb}
\emph{componer} is inherited from Latin \emph{compōnere} \textquotedblleft{}to put together, to arrange\textquotedblright{} (EtymDB), built on \emph{pōnere} \textquotedblleft{}to put\textquotedblright{}.
\end{cousinweb}

\begin{grammarlens}[title={What you've built}]
One more word for telling, explaining and describing at B1.
\end{grammarlens}

\subsection*{Your turn}
\begin{itemize}
\raggedright
  \item \emph{Say it:} \emph{componer}
  \item \emph{Say it:} \emph{componer}, once more
  \item \emph{Say it:} a sentence of your own with \emph{componer}
  \item \emph{Recall:} say the Spanish for the song, then the Spanish for rhythm, then say \emph{componer} again
\end{itemize}

\begin{tcolorbox}[breakable,colback=teal!4,colframe=teal!35!black,title={Before you move on}]
What does componer mean? (\textquotedblleft{}To compose, to write (music)\textquotedblright{}.) Say it once more.
\end{tcolorbox}
